\documentclass[11pt]{article}
\usepackage[T1]{fontenc}
\usepackage{lmodern}
\usepackage[margin=1in,top=.8in,bottom=.85in]{geometry}
\usepackage{amsmath,amssymb,amsthm,booktabs,microtype}
\usepackage[hidelinks]{hyperref}
\newtheorem{theorem}{Theorem}
\newtheorem{lemma}[theorem]{Lemma}
\newtheorem{corollary}[theorem]{Corollary}
\DeclareMathOperator{\pos}{pos}
\title{Recurrence and position distributions in arbitrarily ordered throwback queues}
\author{Blake Stone mathematics repository\\[-2pt]\small Proof developed with OpenAI Codex}
\date{November 26, 2026}
\begin{document}
\maketitle
\begin{abstract}
In a throwback queue the head has value $d$ and moves to position $d$, counting from zero.
For every ordering of distinct positive labels we prove that the leader word is recurrent
if and only if each label $x$ initially occupies a position at most $x$. Every recurrent
leader word is uniformly recurrent. A violation produces an explicit prefix that occurs
exactly once. For arbitrary starting orders every finite block has an ultimately periodic
occurrence indicator. We give an exact formula for the limiting distribution of each
label's position, including the dependence on its starting orbit.
\end{abstract}

\section{Which starting orders give recurrence?}
Let $Q$ be an infinite list of distinct positive integers. Write $u_t$ for its head at time
$t$, and $p_t(x)$ for the position of a label $x$. Moving $u_t=d$ past the next $d$ entries
takes the old position $p$ of any label to
\[
 F_d(p)=\begin{cases}d&p=0,\\p-1&1\le p\le d,\\p&p>d.\end{cases}
\]
A word is recurrent if every finite block that occurs once occurs infinitely often. It is
uniformly recurrent if for each such block the gaps between occurrences are bounded.

\begin{theorem}\label{th:rec}
The following are equivalent: the word $u_0u_1\cdots$ is recurrent, it is uniformly
recurrent, and $p_0(x)\le x$ for every label $x$. If $p_0(x)>x$, the initial word ending
at the first appearance of $x$ has exactly one occurrence.
\end{theorem}

Darnell and Dribus proved individual-label recurrence for all distinct positive queues and
uniform recurrence for increasing inputs \cite[Theorems 2 and 7]{dd}. An increasing queue
satisfies this criterion. The criterion also applies, for example, to $(2,1,3,4,\ldots)$.
The proof of the converse recovers physical positions from an observed finite word.

\begin{lemma}\label{lem:hit}
From every time, each label eventually reaches the head. After reaching the head, $x$ stays
in positions $0,\ldots,x$.
\end{lemma}
\begin{proof}
For completeness, induct strongly on the present position $p$ of a label $x$, simultaneously
over all starting times. If $p>0$, some label $y$ in the first $p$ slots is at least $p$,
since only $p-1$ positive integers are smaller. By induction $y$ reaches the head. Until
$x$ itself reaches the head its position is nonincreasing. If $x$ has still not appeared
when $y$ is thrown, its position decreases to at most $p-1$. Induction now gives the desired
visit of $x$. Position zero is the base case. After a throw of $x$ its position is $x$,
and every other throw either preserves or decreases it.
\end{proof}

\begin{proof}[Proof of Theorem~\ref{th:rec}]
Suppose $p_0(x)>x$ and let $k$ be the first time with $u_k=x$.
Given the head $d$ of a step, the inverse of $F_d$ sends $d$ to $0$, sends $q<d$ to $q+1$,
and fixes $q>d$. In any occurrence of $u_0\cdots u_k$ the final head $x$ is in slot zero.
Inverting the first $k$ specified steps recovers the same starting slot $p_0(x)>x$.
This excludes any start after time $k$ by Lemma~\ref{lem:hit}. A start $s$ with
$1\le s\le k$ would put the actual $x$ at time $k$ at offset $k-s$ in the repeated block,
before its first $x$. Thus only the original occurrence exists.

Now suppose every initial position is within range. Fix a cutoff $M$ and track all labels
$x\le M$, recording their positions. These positions stay in range. The state determines
the next state: a tracked head has its specified throw distance, and an untracked head
moves every tracked label left by one. The new state also determines its predecessor.
If a tracked head $d$ was thrown, it is now in slot $d$. Every smaller tracked label was
crossed and now lies strictly below its own value. Thus the smallest tracked label at its
own value identifies $d$. If there is no such label, the old head was untracked and every
old tracked position is recovered by adding one. For a tracked head, use $F_d^{-1}$.
The forward trajectory therefore has no transient: a collision among finitely many states
can be followed back to time zero. The tracked state and the leader word with larger
values replaced by one common symbol are periodic from time zero. Taking $M$ at least
as large as every value in a specified block makes its occurrences periodic, hence
syndetic. Uniform recurrence implies recurrence.
\end{proof}

\begin{corollary}\label{cor:blocks}
For any starting order, the indicator of the occurrences of each finite nonempty word
is ultimately periodic. Its limiting frequency is rational. It is positive exactly when
the word occurs infinitely often.
\end{corollary}
\begin{proof}
Take the maximum label in the word as cutoff. By Lemma~\ref{lem:hit}, after a common finite
time all tracked labels are within range. The forward and inverse arguments above use only
the bounds for tracked labels. The capped leader word is periodic from that time, and a
literal block with all values below or at the cutoff is recognized from this word.
A periodic indicator with $b$ hits per period $P$ has mean $b/P$.
\end{proof}

\section{Exact position distributions}
Section 7 of Darnell and Dribus \cite{dd} asks how to quantify average positions for
arbitrary distinct inputs. Our earlier note \cite{mix} computes these means for perfectly
mixing increasing prefixes. A prefix here means the \emph{smallest labels}, whatever their
initial locations. The next formula uses only the orbit of the smaller labels and a
congruence. In particular it applies to each orbit separately.

Write the labels in numerical order as $x_0<x_1<\cdots$ and let $r_i=x_i-i+1\ge2$.
To compute a distribution through level $i$, choose a time $T$ by which all labels through
$x_i$ lie in range and stay there, as guaranteed by the lemma. In the queue at time $T+t$,
let $z_i(t)$ be the zero-based position of $x_i$ after deleting all smaller labels.
Equivalently, $x_i$ occupies the $z_i(t)$-th free slot among $0,\dots,x_i$ after the
$i$ smaller labels have been placed. Its possible ranks are $0,\dots,r_i-1$.
Placing labels in order gives a bijection from the product of the rank sets to the
admissible physical configurations.

Call a step \emph{active at level $i$} if its head is at least $x_i$.
Let $P$ be the least period of the smaller configuration, $H$ the number of active steps
in one such period, and $A(t)$ their number in steps $0,\ldots,t-1$, relative to $T$.
For an empty smaller configuration use $P=H=1$. Let $F_t$ be the free slots at time $t$
among $0,\ldots,x_i$, and $\rho_t(v)=\#\{w\in F_t:w<v\}$.

\begin{theorem}\label{th:dist}
Put $g=\gcd(r_i,H)$ and $z=z_i(0)$. The limiting probability of finding $x_i$ in slot $v$
is zero outside $[0,x_i]$ and otherwise equals
\[
 \pi_i(v)=\frac{g}{P r_i}\sum_{t=0}^{P-1}
 {\bf1}_{v\in F_t}\,{\bf1}_{\rho_t(v)\equiv z-A(t)\pmod g}.
\]
Its limiting average position is $\sum_{v=0}^{x_i}v\pi_i(v)$.
The least period and number of active steps for the \emph{next} level are respectively
\[
 P'=P\,r_i/g,\qquad H'=(H/g)(r_i-1).
\]
In particular every admissible orbit through the same set of labels has the same length.
There are $\prod_{j=0}^i r_j/P'$ such orbits.
\end{theorem}
\begin{proof}
A smaller head leaves the relative order of all surviving labels unchanged.
When a head $d\ge x_i$ moves, every smaller label lies before its insertion slot, since
each occupies a position strictly below $x_i$. After deletion of the smaller labels the
throw distance is therefore $d-i$. If $d=x_i$, rank zero becomes $r_i-1$. Otherwise the
head crosses $x_i$ and its rank decreases by one. Thus each active step decrements the
rank cyclically, and $z_i(t)\equiv z-A(t)\pmod{r_i}$. The smaller projection evolves
independently, and its phase determines whether a step is active.

After $P$ steps the smaller configuration returns and the rank is translated by $-H$.
Its order is $r_i/g$. At each base time $0\le t<P$, the visits at $t+kP$ consequently
sample each rank congruent to $z-A(t)$ modulo $g$ exactly once per joint orbit.
On multiplying the uniform time weight $1/P'$ by the indicated count one obtains the
displayed probabilities. The total number of active steps in $P'$ moves is $Hr_i/g$.
Every $r_i$-th active step has leader $x_i$, so $(H/g)(r_i-1)$ pass to the next level.
These arguments apply for every initial rank. The period induction starts with the empty
configuration. The product of the rank-set sizes counts all admissible states, so division
by their common period counts the orbits.
\end{proof}

\paragraph{Dependence on order.}
Take smallest labels $2$ and $4$. The lower period is $3$ with two active steps, so
$r_1=4$, $g=2$, and the joint period is $6$. Starting the queue with $(2,4,5,6,\ldots)$,
the slot counts of $4$ in positions $0,\ldots,4$ are $(1,1,1,2,1)$ per period, giving
mean $13/6$. Starting with $(2,5,4,6,\ldots)$ gives $(1,1,1,1,2)$ and mean $7/3$.
The label set and its limiting head frequencies thus allow different mean positions.

\paragraph{Verification.}
The accompanying Lean development proves the first-visit lemma, the unique-prefix
obstruction, and the two equivalences of Theorem~\ref{th:rec} from the physical position
equations, head identity, positivity, and occupancy of every slot. The distribution formula
is proved above. The exact-arithmetic verifier compares its probabilities with literal
finite projected orbits from every admissible state of small label sets, and separately
runs actual permuted queues.

\begin{thebibliography}{9}
\bibitem{dd} J. Darnell and B. F. Dribus, \emph{Throwback sequences of positive integers},
Journal of Integer Sequences \textbf{28} (2025), Article 25.5.1.
\url{https://cs.uwaterloo.ca/journals/JIS/VOL28/Dribus/dribus4.html}.
\bibitem{mix} B. Stone mathematics repository, \emph{Perfect mixing in increasing
throwback queues} (2026).
\href{https://github.com/blakestone-x/math/tree/6fdc6b705f6c8d04bc1a983871fc10a2afdce3d6/Throwback-mixing}{GitHub manuscript at commit 6fdc6b7}.
\end{thebibliography}
\end{document}
